% ===========================================================================
% Edge-Cloud Collaborative Memory Management
% Usage: \input{section_memory_en} in your main document
% Required packages (ensure loaded in preamble):
%   \usepackage{booktabs, graphicx, amssymb, amsmath}
%   \usepackage{tikz}
%   \usetikzlibrary{arrows.meta, positioning}
% ===========================================================================

% \subsection{Edge-Cloud Collaborative Memory Management}
% \label{sec:memory-sync}

% \subsubsection{Problem Statement}

% During patrol, an on-device robot continuously acquires visual memories, each corresponding to a detected entity in the scene.
% These memories exhibit two inherent dimensions of heterogeneity:

% \begin{enumerate}
%     \item \textbf{Task relevance varies.} Traffic cones and construction barriers should trigger rerouting alerts for other devices; fire hydrants serve as localization landmarks. An undifferentiated memory store cannot support task-aware retrieval.
%     \item \textbf{Privacy sensitivity varies.} Faces and personal belongings are bound to specific individuals---disseminating them without sanitization violates privacy principles. In contrast, environmental information (road hazards, construction zones) is a public resource suitable for sharing.
% \end{enumerate}

% The core challenge is to automatically identify both dimensions at ingest time and enforce a ``never leak, rather under-share'' policy during cross-device synchronization.

% % ---------------------------------------------------------------------------
% \subsubsection{Method}
% \label{subsec:method}

% \paragraph{Two-Dimensional Classification}
% Rather than deferring relevance and privacy decisions to computationally expensive vision-language models, we adopt a prior-knowledge-driven approach: a compact set of task-specific priors defines which entities merit storage (task relevance) and, among stored records, which may be shared across devices (privacy sensitivity).
% This serves as a lightweight selective gating mechanism that eliminates model inference from the critical path while remaining extensible to model-assisted refinement.
% Each memory record is accordingly annotated with two orthogonal attributes (Table~\ref{tab:dimensions}), determined by deterministic rules and keyword matching.

% \begin{table}[ht]
% \centering
% \caption{Two-dimensional memory classification schema.}
% \label{tab:dimensions}
% \begin{tabular}{lll}
% \toprule
% Dimension & Values & Determination \\
% \midrule
% \texttt{sensitivity} & high / low & Whether bound to a specific individual \\
% \texttt{nav\_category} & hazard / poi / landmark / person / object & Task relevance \\
% \bottomrule
% \end{tabular}
% \end{table}

% A record is shareable if and only if its sensitivity is low; only shareable records are eligible for promotion to common memory.

% The open-vocabulary detector is configured with a task-specific keyword list that encodes task priors; only entities matching predefined categories (e.g., traffic cone $\mapsto$ hazard, fire hydrant $\mapsto$ landmark) are detected and enter the memory pipeline---front-gating irrelevant observations at perception time.
% Because the detector operates on natural-language prompts, novel entity types can be incorporated by extending the keyword list without retraining or rule modification, yielding a degree of generalization within the prior-based framework.

% The current instantiation applies three priority-ordered \textbf{sensitivity rules}.
% A record is classified as \textit{high} sensitivity if any of the following conditions holds; otherwise it defaults to \textit{low}:
% \begin{enumerate}
%     \item The entity has a \texttt{belongs\_to} relation referencing a person entity;
%     \item The entity category is \texttt{person}, or a face embedding is associated;
%     \item The description text contains a known person name or matches a privacy keyword.
% \end{enumerate}

% % ---------------------------------------------------------------------------
% \paragraph{Private--Common Memory Synchronization}
% The system separates memories into \textit{private memory} (edge-local, per-device) and \textit{common memory} (cloud-hosted, shared across devices).
% As shown in Figure~\ref{fig:sync-arch}, synchronization follows an offline-first protocol.

% \begin{figure}[t]
% \centering
% \begin{tikzpicture}[
%     node distance=0.8cm and 2.8cm,
%     box/.style={draw, rounded corners, minimum width=3.0cm, minimum height=2.0cm, align=center, font=\small},
%     arrow/.style={-{Stealth[length=3mm]}, thick},
%     lbl/.style={font=\footnotesize, align=center}
% ]
%     % Nodes
%     \node[box] (A) {Private Memory\\(Device A)\\[3pt]\footnotesize detect $\rightarrow$ classify\\[-2pt]\footnotesize $\rightarrow$ upload};
%     \node[box, right=of A] (cloud) {Common Memory\\(Cloud)\\[3pt]\footnotesize store \&\\[-2pt]\footnotesize deduplicate};
%     \node[box, right=of cloud] (B) {Private Memory\\(Device B)\\[3pt]\footnotesize download $\rightarrow$ dedup\\[-2pt]\footnotesize $\rightarrow$ merge};

%     % Arrows
%     \draw[arrow] (A) -- node[above, lbl] {upload\\(shareable only)} (cloud);
%     \draw[arrow] (cloud) -- node[above, lbl] {query \& fetch} (B);
% \end{tikzpicture}
% \caption{Overview of private--common memory synchronization. Only low-sensitivity records are promoted to common memory; downloaded records are marked non-re-shareable to prevent uncontrolled propagation.}
% \label{fig:sync-arch}
% \end{figure}

% Three design principles govern the architecture:
% \begin{itemize}
%     \item \textbf{Offline-first:} Private memory is the single source of truth; full functionality is preserved without network connectivity.
%     \item \textbf{Private by default:} Memories remain in private storage unless explicitly classified as low-sensitivity and promoted to common memory.
%     \item \textbf{No re-propagation:} Records downloaded from common memory are automatically marked as non-shareable.
% \end{itemize}

% % ---------------------------------------------------------------------------
% \paragraph{Receiver-Side Deduplication}
% The receiving device applies a two-strategy deduplication check (embedding cosine similarity and keyword Jaccard overlap, with configurable thresholds) on each incoming record to suppress redundant entries.
% Records surviving both checks are merged into private memory as non-shareable.

% % ---------------------------------------------------------------------------
% \subsubsection{Functional Verification}

% We verify the pipeline logic on two synthetic scenarios (4 frames each, generated via text-to-image synthesis) using an open-vocabulary detector for entity extraction.
% A total of 29 memory records are extracted across 7 entity categories.
% All records are classified consistently with the rule design: the 5 person-category records are correctly assigned high sensitivity and blocked from promotion to common memory, while all 24 environmental records receive low-sensitivity labels and pass the shareability gate.
% Dynamic vocabulary extension enables detection of user-specified categories beyond the default set (e.g., bench, thermos), confirming that the prior-based classification accommodates novel entity types without rule modification.

% % ---------------------------------------------------------------------------
% \subsubsection{Summary}

% This module implements automatic two-dimensional classification and privacy-safe cloud synchronization of visual memories.
% The classified memory store provides a standardized query interface (category-filtered, location-aware) designed to support downstream navigation and planning modules with task-specific retrieval without exposing raw sensor data.


\subsection{Edge--Cloud Collaborative Memory Management}
\label{sec:memory-sync}

Embodied agents continuously acquire heterogeneous memories during interaction with the physical world, including maps, semantic landmarks, obstacles, objects, and human-related observations. 
While public environmental memories, such as roadblocks, construction areas, and navigational landmarks, can improve multi-agent collaboration when shared, privacy-sensitive memories, such as faces, personal belongings, or person-associated objects, must remain local. 
To address this conflict, we design an edge--cloud collaborative memory system that separates private memory on the robot from common memory in the cloud.

\paragraph{Memory Partition.}
Each robot maintains a \emph{private memory} at the edge, which stores all locally perceived memories, including map memory, semantic memory, and multi-modal memory. 
The cloud maintains a \emph{common memory} that only contains public, low-sensitivity memories useful for collective navigation and task execution. 
This design follows a private-by-default principle: a memory item is never uploaded unless it is explicitly classified as shareable.

\paragraph{Privacy-aware Memory Gating.}
For each newly generated memory item, the system performs a privacy judgment before synchronization. 
A memory is regarded as non-shareable if it contains or is associated with personally identifiable information, including persons, faces, names, personal objects, or ownership relations to a specific individual. 
In contrast, public environmental information, such as maps, traffic cones, barriers, road damage, and static landmarks, is classified as shareable.


To evaluate the reliability of the privacy protection mechanism, we constructed a dedicated dataset for privacy classification and upload-decision testing. Experimental results show that the system achieves over 99\% accuracy in identifying whether a memory item is suitable for cloud synchronization. This high-accuracy privacy judgment significantly reduces the risk of uploading sensitive user or environment-specific information, while still allowing non-sensitive public knowledge to be shared across robots. As a result, the proposed memory system balances two key requirements of embodied intelligence: preserving user privacy through edge-side private memory, and enhancing group-level collaboration through cloud-side common memory.

